\documentclass[conference, 10pt]{IEEEtran}

\usepackage{cite}
\usepackage{amsmath,amssymb,amsfonts}
\usepackage{graphicx}
\usepackage{tabularx}
\usepackage{parskip}
\usepackage{float}
\usepackage{booktabs}
\usepackage{hyperref}

\begin{document}

\title{Evaluating Bagging Ensembles for Cyberattack Classification}

\author{
\IEEEauthorblockN{AM. Zenzile}
\IEEEauthorblockA{Stellenbosch University}
\IEEEauthorblockA{27403122@sun.ac.za}
}

\maketitle

\begin{abstract}
    Ensemble learning is based on the principle that combining multiple learning models can produce better 
    predictions than relying on a single model. This assignment investigates the effect of ensemble learning 
    on cyberattack classification, focusing on two questions: whether combining multiple instances of the same 
    base classifier improves classification performance compared with an individual classifier, and how the choice 
    of sampling strategy affects ensemble performance. The investigation compares individual K-nearest neighbours 
    (KNN) and decision tree classifiers with bagging ensembles comprising multiple instances of the respective 
    classifiers. The effect of sampling strategy is then examined by comparing bootstrap sampling with sampling 
    without replacement. Results indicate that bagging does not necessarily improve classification performance by 
    a huge margin, with its benefits depending on the base algorithm, while the choice 
    of sampling strategy had only a marginal effect.
\end{abstract}

\begin{IEEEkeywords}
classification tree, K-nearest neighbours, machine learning, network traffic classification
\end{IEEEkeywords}

\section{Introduction}

Network traffic classification is a supervised machine learning problem in which network connections are assigned 
to predefined attack categories. Classification algorithms such as K-nearest neighbours (KNN) and decision trees 
provide approaches to solving this problem. Ensemble learning extends these approaches by combining predictions 
from multiple models to produce a single prediction. Bagging, or bootstrap aggregating, is an ensemble method 
that trains multiple instances of a base classifier on different samples of the training data and aggregates 
their predictions.

Although bagging combines multiple classifiers to improve predictive performance, its effectiveness depends on 
how the ensemble is constructed. The choice of sampling strategy determines the training data presented to each 
base classifier and may influence the diversity of the resulting ensemble. Bootstrap sampling draws observations 
with replacement, whereas sampling without replacement selects observations without duplication within each sample. 
The resulting differences in training samples raise questions about the effect of sampling strategy on classification 
performance. Bagging also does not guarantee an improvement over an individual classifier, and the extent of any 
improvement may differ between learning algorithms.

The aim of this assignment therefore, is to evaluate bagging for cyberattack classification by addressing two questions: (1) 
Does bagging improve classification performance compared with individual classifiers, and does the effect differ 
between K-nearest neighbours (KNN) and decision trees? (2) How does bootstrap sampling compare with sampling 
without replacement in terms of ensemble classification performance? The first question is investigated by 
comparing individual KNN and decision tree classifiers with bagging ensembles built from the respective base 
classifiers. The findings suggest that combining multiple classifiers does not necessarily improve classification 
performance, as the benefits of bagging depend on the base algorithm and remain marginal in this experiment. 
Although subsampling without replacement generally outperformed bootstrap sampling, the small differences 
indicate that the choice of sampling strategy also has limited effect on overall performance.

\section{Background}

\subsection{k-Nearest Neighbours}

k-Nearest Neighbours (kNN) is a supervised, non-parametric, similarity-based learning approach. 
The method retains the labelled training instances and does not construct an explicit parametric 
model during training \cite{Cover1967}. The parameter (k) specifies the number of nearest training 
instances considered when classifying a new instance. Classification is performed by identifying the 
(k) training instances most similar to the new instance and assigning the class supported by majority of the 
neighbours. Consequently, the value of (k) influences the inductive bias of the classifier. 
A small value of (k) produces a more locally sensitive classifier, while a large value of (k) 
produces a smoother classification rule. The choice of (k) therefore affects the balance between 
sensitivity to individual training instances and generalisation across the dataset \cite{GarciaPedrajas2017}.

The notion of similarity between instances is determined by a distance or similarity measure. 
The choice of measure depends on the representation and nature of the features being compared. 
Common measures for numerical data include Euclidean, Manhattan, and Minkowski distance, while 
measures such as Jaccard similarity are used for appropriate binary representations. Different 
distance measures can produce different neighbourhood relationships and therefore affect the 
classification performance of kNN \cite{AbuAlfeilat2019}. During prediction, the feature representation 
of a new instance is compared with the stored training instances, the nearest instances are identified 
according to the selected measure, and the class of the new instance is determined from the 
resulting neighbourhood.

\subsection{Classification Decision Trees}

A classification decision tree also is a supervised, non-parametric learning approach with the 
distinction that it represents a classification policy or the model as a tree structure. The 
tree consists of internal nodes that apply tests to input features and leaf nodes that represent 
predicted classes. A classification decision tree is constructed by recursively partitioning the 
training data  according to feature values. The resulting root-to-leaf paths represent decision 
rules that can be used to classify new instances \cite{Singh2014}.

The inductive bias of a decision tree is influenced by the conditions used to control tree growth. 
Tree depth and stopping conditions, such as the minimum number of instances permitted in a leaf, 
determine the complexity of the resulting classification rules. Less restrictive conditions allow 
the tree to represent increasingly detailed patterns in the training data, while more restrictive 
conditions produce simpler trees. The splitting criterion also influences the structure of the 
resulting tree. Information gain measures the reduction in class uncertainty produced by a split 
and forms the basis of the ID3 approach. Gain ratio extends information gain by accounting for the 
number of possible outcomes of a split and is associated with the C4.5 approach \cite{Singh2014}. 
The Gini index provides another measure of class impurity and is used as a splitting criterion in 
classification trees such as CART \cite{SivagamaSundhari2011}. These criteria determine how candidate 
feature splits are evaluated during tree construction.

\subsection{Bagging as an Ensemble Learning Technique} 
Ensemble learning on the other end, is not a model or a training algorithm but a machine learning approach 
that combines the predictions of multiple learning models to produce a single prediction. The underlying 
principle is that combining the outputs of multiple models can improve predictive performance compared with 
relying on an individual model. An ensemble consists of a collection of models, referred to as base learners, 
whose predictions are combined using an aggregation strategy. Ensemble methods differ in how the base learners 
are constructed and how their predictions are combined. Bagging, or bootstrap aggregating, is one such technique, 
in which multiple instances of a base learning algorithm are trained, independently, on different samples of the 
training data before their predictions are aggregated. \cite{DBLP:conf/mcs/Dietterich00}

Bagging constructs an ensemble by sampling observations from the training dataset to create separate training 
subsets for the base learners. Sampling introduces variation between the training subsets, allowing the base 
learners to learn from different representations of the original data. Each base learner is trained independently 
on its assigned subset, and the resulting models are combined during prediction. For classification, majority voting 
assigns the class predicted by the largest number of base learners. Weighted majority voting assigns weights to the 
base learners' votes, allowing models with higher assigned weights to contribute more to the final prediction. 
The predicted class is the one with the greatest total weight. For regression, predictions can be aggregated using 
the median, which selects the middle value after the individual predictions are ordered. The median is less sensitive 
to extreme predictions than the arithmetic mean. Breiman's original formulation of bagging uses plurality voting for 
classification and averaging for numerical predictions, although other aggregation strategies are possible 
\cite{DBLP:journals/ml/Breiman96b}.

\subsection{Sampling Strategies}
There are two primary sampling strategies widely used during bagging, bootstrap sampling (with replacement) and 
subsampling without replacement. 

\textbf{Bootstrap sampling.} Bootstrap sampling creates a training sample by drawing observations from the original dataset with 
replacement. An observation can therefore appear multiple times, while others are excluded. This introduces variation between the 
training samples used to train the base learners \cite{DBLP:journals/ml/Breiman96b}. However, duplicate observations can give 
particular instances additional influence, including incorrectly labelled instances. This is especially an issue with a dataset 
that has class imbalance or skewed distribution. The effectiveness of bootstrap sampling therefore consequently depends on the dataset's 
characteristics and the learning algorithm.

\textbf{Subsampling without replacement.} Subsampling on the other hand, selects observations without replacement, ensuring that each 
observation appears at most once within a sample. Compared with bootstrap sampling, this avoids duplicate observations and provides more 
distinct training instances at the same sample size. Subsampling can also improve robustness to class-label noise under certain conditions 
\cite{sabzevari2018two}. However, smaller samples can exclude informative or minority-class observations, potentially reducing classification 
performance. Its effectiveness therefore depends on the sample size and the dataset's characteristics as well.

\section{Methodology}

The assignment investigates the effectiveness of bagging ensembles for cyberattack classification through two main research questions:

\textbf{RQ1: Effect of bagging.} Does bagging improve classification performance compared with individual classifiers, and does the effect 
differ between K-nearest neighbours (KNN) and decision trees?

This question is investigated by comparing each individual classifier with a bagging ensemble constructed from the same base learning algorithm. 
Subsampling is used for both ensembles to maintain a consistent sampling strategy during the comparison.

\textbf{RQ2: Effect of sampling strategy.} How does bootstrap sampling compare with sampling without replacement in terms of ensemble 
classification performance?

This question is investigated by comparing KNN and decision tree ensembles constructed using the two sampling strategies. 
The comparison examines whether the choice of sampling strategy affects the classification performance of either ensemble.

\subsection{Dataset}
The cyberattack dataset used contains 257,673 instances and 44 columns, including 43 descriptive features and 
the target feature, \texttt{attack\_cat}. The descriptive features comprise 39 numerical features and four nominal 
features: \texttt{proto}, \texttt{service}, \texttt{state} and \texttt{is\_ftp\_login}. The target represents ten 
attack categories. The dataset contains missing values exclusively in \texttt{service}, where 141,321 instances 
(54.85\%) lack a recorded value. The target classes are also imbalanced, with class proportions ranging from 36.09\% 
for the largest class to 0.07\% for the smallest. A complete description of the dataset features is provided in Appendix 
\ref{sec: dataset}. These characteristics informed the preprocessing procedure applied before model training and evaluation.

\subsection{Pre-Processing} 
The dataset was first divided into training and testing sets using an 80:20 stratified split~\cite{DBLP:conf/ijcai/Kohavi95} 
with a random seed of 42. The \texttt{id} column was removed because it identifies individual records, and does not hold 
valuable information for identifying a cyberattack category. To reduce feature redundancy, Pearson correlation was used to 
identify highly correlated feature pairs, and mutual information with the target variable, \texttt{attack\_cat}, was 
calculated for the features in each pair using scikit-learn's \texttt{mutual\_info\_classif} function. Mutual information 
measures the statistical dependence between a feature and the target, providing a basis for retaining the feature with the 
stronger association with the target. Following that, the \texttt{ct\_ftp\_cmd} feature was removed because its mutual 
information score (0.0053) was lower than that of \texttt{is\_ftp\_login} (0.0080). Similarly, \texttt{sloss} and \texttt{dloss} were 
removed because their scores (0.2914 and 0.3190, respectively) were lower than those of \texttt{sbytes} (1.0886) and 
\texttt{dbytes} (0.5895). The \texttt{service} feature contained missing values in 54.85\% of instances; these were 
replaced with an explicit \texttt{\_\_MISSING\_\_} category to retain the feature without discarding a substantial 
proportion of the dataset.

The remaining features were transformed according to the requirements of each classifier. For KNN, numerical features 
were standardised using \texttt{StandardScaler} from scikit-learn~\cite{pedregosa2011scikit}, fitted only on the training 
data, because KNN relies on distance-based similarity measures that can be dominated by features with larger numerical scales, 
which is much of a case with the considered dataset. Nominal features were one-hot encoded to represent categories without 
imposing an artificial ordering. For decision trees, nominal features were also one-hot encoded, while numerical features 
were left unscaled because tree splits are not sensitive to feature magnitude in the same way as distance-based KNN similarity 
measures. All transformations were fitted using the training data and then applied to the test data without refitting, 
preventing information leakage during evaluation.

\subsection{Emperical Process}
The experiments were conducted in two stages to address the research questions. The first experiment investigated the 
effect of sampling strategy by comparing bootstrap sampling with subsampling without replacement for KNN and decision 
tree ensembles. Each ensemble comprised five base learners, with each learner trained on a sample containing 80\% of the 
training data. The results were compared to identify the better-performing sampling strategy for each ensemble type. 
The second experiment evaluated the effect of bagging by comparing each individual classifier with its corresponding ensemble, 
using the sampling strategy selected from the first experiment. This approach allowed the main experiment to use the 
better-performing sampling configuration identified during the preliminary comparison. Ensemble predictions were determined 
through majority voting, and the same training and testing partitions were used throughout both experiments to ensure consistent 
evaluation.

Classification performance was evaluated using accuracy, macro-precision, macro-recall and macro-F1~\cite{DBLP:journals/ipm/SokolovaL09}. 
Accuracy measures the proportion of correctly classified instances, while precision measures the proportion of 
positive predictions that are correct for a given class. Recall measures the proportion of instances belonging 
to a class that are correctly identified. Macro-F1 calculates the F1-score for each class and takes their 
unweighted mean, giving each class equal importance regardless of its frequency. Macro-precision and macro-recall 
were also calculated as unweighted means across classes to account for the imbalanced target distribution. 
The results were recorded for each configuration and examined using metric comparison tables, heatmaps and 
paired bar charts.

\subsection{Implementation}

The implementation was organised into three main components: preprocessing, model construction and experiment execution. 
The preprocessing component handled the train-test split and prepared separate feature representations for KNN and decision 
trees. The model component, implemented in \texttt{models.py}, used scikit-learn implementations of KNN and decision 
trees~\cite{pedregosa2011scikit} and comprised four classes: \texttt{SamplingStrategies}, which generated bootstrap samples 
and subsamples without replacement; \texttt{IndividualModels}, which trained the baseline classifiers; \texttt{BaggedKNN}, 
which constructed KNN ensembles; and \texttt{BaggedDecisionTree}, which constructed decision tree ensembles. Both ensemble 
classes reused shared sampling, training and majority-voting procedures to maintain a consistent implementation across the 
two algorithms. An experiment runner, \texttt{runner.py}, acted as the interface between these components, coordinating data 
preparation, model training, prediction, evaluation and output generation.

The experiments used fixed hyperparameters rather than tuning the models. The individual KNN classifier and each KNN ensemble 
learner used \(k=20\), while the individual decision tree and each tree in the corresponding ensemble had a maximum depth of 10. 
Each ensemble contained five base learners, trained on separate samples comprising 80\% of the training partition. The sampling 
strategy was configurable, allowing the ensembles to use either bootstrap sampling or subsampling without replacement. A random 
seed of 42 was used to support reproducibility. During prediction, each base learner contributed one class-label vote, and the 
ensemble returned the class receiving the most votes. The runner calculated accuracy, macro-precision, macro-recall and macro-F1 
or each configuration, recorded the predictions and metric results in CSV files, and generated heatmaps and paired bar charts in 
PDF format. Command-line arguments allowed either experiment to be run independently or both to be executed in sequence, with the 
sampling strategy for the main experiment specified by the user. The full implementation can be found in this GitHub 
repository: \url{https://github.com/APZenzile/ML_assignment_4}.

\section{Results}

\subsection{Experiment 1: Effect of Sampling Strategy}

\begin{figure}[H]
  \centering
  \includegraphics[width=\columnwidth]{../output/figures/rq2_heatmap.pdf}
  \caption{Bagged KNN and decision tree performance under bootstrap sampling and subsampling.}
  \label{fig:exp1_heatmap}
\end{figure}

The heatmap shows that subsampling without replacement achieved higher accuracy, macro-precision and macro-recall 
for both ensemble types. For KNN, macro-F1 increased from 0.4200 to 0.4242, whereas for decision trees it decreased 
slightly from 0.4876 to 0.4868. Accuracy changed only marginally for both algorithms. These results suggest that 
subsampling without replacement offers a small performance advantage overall, although its effect on macro-F1 
depends on the base algorithm.

\begin{figure}[H]
    \centering
    \includegraphics[width=\columnwidth]{../output/figures/rq2_decision_tree_sampling_comparison.pdf}
    \caption{Decision Tree performance with different sampling techniques.}
    \label{fig:exp1_dt_sampling}
\end{figure}

For the decision tree ensemble, subsampling without replacement produced marginally higher accuracy 
(0.7989 versus 0.7987), macro-precision (0.7251 versus 0.6890) and macro-recall (0.4760 versus 0.4701) 
than bootstrap sampling. However, bootstrap sampling achieved a slightly higher macro-F1 (0.4876 versus 0.4868). 
These results suggest that subsampling changes the balance between precision and recall but does not 
consistently improve all evaluation metrics.

\begin{figure}[H]
    \centering
    \includegraphics[width=\columnwidth]{../output/figures/rq2_knn_sampling_comparison.pdf}
    \caption{KNN performance on different sampling techniques.}
    \label{fig:exp1_knn_sampling}
\end{figure}

For the KNN ensemble, subsampling without replacement outperformed bootstrap sampling across all four metrics, 
although the differences were small. Accuracy increased from 0.7679 to 0.7696, while macro-F1 increased from 
0.4200 to 0.4242. This suggests that avoiding duplicate observations within each training sample provides a 
marginal benefit for KNN in this experiment. However, the limited improvement indicates that the sampling 
strategy had little influence on the ensemble's overall performance.

\subsection{Experiment 2: Individual Models versus Bagged Ensembles}

\begin{figure}[H]
    \centering
    \includegraphics[width=\columnwidth]{../output/figures/rq1_heatmap.pdf}
    \caption{Bagged KNN and Classification Tree against individual respective models.}
    \label{fig:exp2_heatmap}
\end{figure}

The heatmap shows that bagging did not consistently improve classification performance across the two base algorithms. 
The KNN ensemble achieved slightly higher accuracy than individual KNN (0.7696 versus 0.7693), but its macro-F1 
decreased from 0.4270 to 0.4242. In contrast, the decision tree ensemble achieved marginally higher accuracy, 
macro-recall and macro-F1 than the individual tree, with macro-F1 increasing from 0.4851 to 0.4868. These 
results suggest that the benefit of bagging depends on the base algorithm and that combining multiple classifiers 
does not necessarily produce a meaningful performance improvement.

\begin{figure}[H]
    \centering
    \includegraphics[width=\columnwidth]{../output/figures/rq1_knn_comparison.pdf}
    \caption{KNN ensemble vs individual KNN.}
    \label{fig:exp2_knn_comp}
\end{figure}

The individual KNN classifier outperformed the bagged KNN ensemble in macro-precision (0.6227 versus 0.5706), 
macro-recall (0.4098 versus 0.4070) and macro-F1 (0.4270 versus 0.4242). The ensemble achieved marginally higher 
accuracy, increasing from 0.7693 to 0.7696. The reductions in macro-precision and macro-F1 suggest that bagging 
offered no overall benefit to KNN under the selected configuration, despite the slight increase in accuracy.

\begin{figure}[H]
    \centering
    \includegraphics[width=\columnwidth]{../output/figures/rq1_decision_tree_comparison.pdf}
    \caption{Classification tree ensemble vs individual classification tree.}
    \label{fig:exp2_dt_comp}
\end{figure}

Decision tree ensemble achieved slightly higher accuracy (0.7989 versus 0.7988), macro-precision (0.7251 versus 0.7191), 
macro-recall (0.4760 versus 0.4678) and macro-F1 (0.4868 versus 0.4851) than the individual decision tree. Although the 
ensemble improved all four metrics, the differences were small. This suggests that bagging provided a modest improvement 
for decision trees, rather than the substantial performance gain that might be expected from combining multiple classifiers.

\appendices
  
\subsection{Dataset}
\label{sec: dataset}

% ---------- Part 1 ----------
\begin{table}[H]
    \caption{Dataset features (part 1 of 2).}\label{tab:features1}
    \scriptsize
    \setlength{\tabcolsep}{3pt}
    \renewcommand{\arraystretch}{1.1}
    \begin{tabularx}{\columnwidth}{@{}>{\ttfamily}l l >{\raggedright\arraybackslash}X@{}}
    \toprule
    \textbf{Name} & \textbf{Type} & \textbf{Description} \\
    \midrule
    proto & Nominal & Transaction protocol \\ \addlinespace[4pt]
    state & Nominal & Indicates to the state and its dependent protocol, e.g. ACC, CLO, CON and (-) (if not used state) \\ \addlinespace[4pt]
    dur & Float & Record total duration \\ \addlinespace[4pt]
    sbytes & Integer & Source to destination transaction bytes \\ \addlinespace[4pt]
    dbytes & Integer & Destination to source transaction bytes \\ \addlinespace[4pt]
    sttl & Integer & Source to destination time to live value \\ \addlinespace[4pt]
    dttl & Integer & Destination to source time to live value \\ \addlinespace[4pt]
    sloss & Integer & Source packets retransmitted or dropped \\ \addlinespace[4pt]
    dloss & Integer & Destination packets retransmitted or dropped \\ \addlinespace[4pt]
    service & Nominal & http, ftp, smtp, ssh, dns, ftp-data, irc and (-) if not much used service \\ \addlinespace[4pt]
    Sload & Float & Source bits per second \\ \addlinespace[4pt]
    Dload & Float & Destination bits per second \\ \addlinespace[4pt]
    Spkts & Integer & Source to destination packet count \\ \addlinespace[4pt]
    Dpkts & Integer & Destination to source packet count \\ \addlinespace[4pt]
    swin & Integer & Source TCP window advertisement value \\ \addlinespace[4pt]
    dwin & Integer & Destination TCP window advertisement value \\ \addlinespace[4pt]
    stcpb & Integer & Source TCP base sequence number \\ \addlinespace[4pt]
    dtcpb & Integer & Destination TCP base sequence number \\ \addlinespace[4pt]
    smean & Integer & Mean of the flow packet size transmitted by the src \\ \addlinespace[4pt]
    dmean & Integer & Mean of the flow packet size transmitted by the dst \\ \addlinespace[4pt]
    trans\_depth & Integer & Represents the pipelined depth into the connection of http request/response transaction \\ \addlinespace[4pt]
    response\_body\_len & Integer & Actual uncompressed content size of the data transferred from the servers http service. \\
    \bottomrule
    \end{tabularx}
\end{table}
    
% ---------- Part 2 ----------
\begin{table}[H]
    \caption{Dataset features (part 2 of 2).}\label{tab:features2}
    \scriptsize
    \setlength{\tabcolsep}{3pt}
    \renewcommand{\arraystretch}{1.1}
    \begin{tabularx}{\columnwidth}{@{}>{\ttfamily}l l >{\raggedright\arraybackslash}X@{}}
    \toprule
    \textbf{Name} & \textbf{Type} & \textbf{Description} \\
    \midrule
    Sjit & Float & Source jitter (mSec) \\ \addlinespace[4pt]
    Djit & Float & Destination jitter (mSec) \\ \addlinespace[4pt]
    sinpkt & Float & Source interpacket arrival time (mSec) \\ \addlinespace[4pt]
    dinpkt & Float & Destination interpacket arrival time (mSec) \\ \addlinespace[4pt]
    tcprtt & Float & TCP connection setup round-trip time, the sum of synack and ackdat. \\ \addlinespace[4pt]
    synack & Float & TCP connection setup time, the time between the SYN and the SYN\_ACK packets. \\ \addlinespace[4pt]
    ackdat & Float & TCP connection setup time, the time between the SYN\_ACK and the ACK packets. \\ \addlinespace[4pt]
    is\_sm\_ips\_ports & Binary & If source (1) and destination (3) IP addresses equal and port numbers (2)(4) equal then, this variable takes value 1 else 0 \\ \addlinespace[4pt]
    ct\_state\_ttl & Integer & No. for each state (6) according to specific range of values for source/destination time to live (10) (11). \\ \addlinespace[4pt]
    ct\_flw\_http\_mthd & Integer & No. of flows that has methods such as Get and Post in http service. \\ \addlinespace[4pt]
    is\_ftp\_login & Binary & If the ftp session is accessed by user and password then 1 else 0. \\ \addlinespace[4pt]
    ct\_ftp\_cmd & Integer & No of flows that has a command in ftp session. \\ \addlinespace[4pt]
    ct\_srv\_src & Integer & No. of connections that contain the same service (14) and source address (1) in 100 connections according to the last time (26). \\ \addlinespace[4pt]
    ct\_srv\_dst & Integer & No. of connections that contain the same service (14) and destination address (3) in 100 connections according to the last time (26). \\ \addlinespace[4pt]
    ct\_dst\_ltm & Integer & No. of connections of the same destination address (3) in 100 connections according to the last time (26). \\ \addlinespace[4pt]
    ct\_src\_ltm & Integer & No. of connections of the same source address (1) in 100 connections according to the last time (26). \\ \addlinespace[4pt]
    ct\_src\_dport\_ltm & Integer & No of connections of the same source address (1) and the destination port (4) in 100 connections according to the last time (26). \\ \addlinespace[4pt]
    ct\_dst\_sport\_ltm & Integer & No of connections of the same destination address (3) and the source port (2) in 100 connections according to the last time (26). \\ \addlinespace[4pt]
    ct\_dst\_src\_ltm & Integer & No of connections of the same source (1) and the destination (3) address in 100 connections according to the last time (26). \\ \addlinespace[4pt]
    attack\_cat & Nominal & The name of each attack category. In this data set, nine categories e.g. Fuzzers, Analysis, Backdoors, DoS, Exploits, Generic, Reconnaissance, Shellcode and Worms \\
    \bottomrule
    \end{tabularx}
\end{table}

\bibliographystyle{IEEEtran}
\bibliography{references}

\end{document}